% !TEX root = 0000.tex
% Anexo D
% Contenido tomado de main_2.tex (líneas 3882-3967) sin cambios de texto.

% El texto está en inglés: reglas de corte de palabra del inglés
\begin{hyphenrules}{english}

\chapter{Anexo: Tasks of the PINNacle Benchmark}\chaptermark{Tasks of the PINNacle Benchmark}\label{app:benchmark}

This appendix lists the twelve tasks of Hao et al.~\cite{hao2024pinnacle} used in Section~\ref{sec:exp-benchmark}
(CI: initial condition; CF: boundary condition).

\section{Burgers Equation}
\paragraph{Burgers1d-C.} $u_t+uu_x=\nu u_{xx}$, $(x,t)\in[-1,1]\times[0,1]$, $u(x,0)=-\sin(\pi x)$,
$u(-1,t)=u(1,t)=0$, $\nu=0.01/\pi$.

\paragraph{Burgers2d-C (coupled).} $\mathbf{u}_t+\mathbf{u}\cdot\nabla\mathbf{u}-\nu\Delta\mathbf{u}=0$,
$(x,y,t)\in[0,L]^2\times[0,1]$, $\mathbf{u}(x,y,0)=2\mathbf{w}(x,y)+\mathbf{c}$ with $\mathbf{w}$ a random
Fourier series (Gaussian random field), periodic boundary conditions, $L=4$, $\nu=0.001$.

\section{Poisson Equation (elliptic)}
\paragraph{Poisson2d-C.} $-\Delta u=0$ in $\Omega=[-0.5,0.5]^2\setminus\bigcup_{i=1}^{4}R_i$, where $R_i$
are four circles of radius $0.1$ centered at $(\pm0.3,\pm0.3)$; $u=0$ on $\partial R_i$ and $u=1$ on the
outer rectangle.

\paragraph{Poisson2d-CG (Poisson--Boltzmann).} $-\Delta u+k^2u=f(x,y)$ in
$[-1,1]^2\setminus\bigcup_{i=1}^{4}R_i$, $f(x,y)=A\sum_i(\mu_i^2+x_i^2)\sin(\mu_1\pi x)\sin(\mu_2\pi y)$;
$u=0.2$ on the rectangle and $u=1$ on the circles; $\mu_1=1$, $\mu_2=4$, $k=8$, $A=10$.

\paragraph{Poisson2d-MS (multi-scale).} $-\nabla\cdot(a(\mathbf{x})\nabla u)=f(x,y)$ in $[-10,10]^2$,
$\partial u/\partial\mathbf{n}+u=0$ on $\partial\Omega$, with $a(\mathbf{x})$ piecewise linear on square subdomains.

\section{Heat Equation (parabolic)}
\paragraph{Heat2d-CG (complex geometry).} $u_t-\Delta u=0$ in
$\big([-8,8]\times[-12,12]\setminus\bigcup_iR_i\big)\times[0,3]$, Robin condition
$-\mathbf{n}\cdot(-c\nabla u)=g-qu$ with $c=1$; large circles centered at $(\pm4,\pm3)$, $(\pm4,\pm9)$, $(0,0)$,
$(0,\pm6)$ with $r=1$, $g=5$, $q=1$; small circles centered at $(\pm3.2,\pm6)$, $(\pm3.2,0)$ with $r=0.4$,
$g=1$, $q=1$.

\paragraph{Heat2d-MS (multi-scale).} $u_t-\frac{1}{(500\pi)^2}u_{xx}-\frac{1}{\pi^2}u_{yy}=0$ in
$[0,1]^2\times[0,5]$, $u(x,y,0)=\sin(20\pi x)\sin(\pi y)$, $u=0$ on $\partial\Omega$.

\paragraph{Heat2d-VC (variable coefficients).} $u_t-\nabla\cdot(a(\mathbf{x})\nabla u)=f(\mathbf{x},t)$ in
$[0,1]^2\times[0,5]$, $f=A\sin(m_1\pi x)\sin(m_2\pi y)\sin(m_3\pi t)$ with $A=200$, $m_1=1$, $m_2=5$,
$m_3=1$; $a(\mathbf{x})$ generated by an exponential Gaussian random field; $u(x,y,0)=0$, $u=0$ on $\partial\Omega$.

\section{Navier--Stokes Equations}
\paragraph{NS2d-C (lid-driven cavity).} $\mathbf{u}\cdot\nabla\mathbf{u}+\nabla p-\frac{1}{Re}\Delta\mathbf{u}=0$,
$\nabla\cdot\mathbf{u}=0$ in $[0,1]^2$; on the top lid $\mathbf{u}=(4x(1-x),0)$; elsewhere $\mathbf{u}=(0,0)$,
and $p=0$ at $(0,0)$; $Re=100$.

\paragraph{NS2d-CG (backward-facing step).} Same equations in
$\Omega=[0,4]\times[0,2]\setminus\big([0,2]\times[1,2]\big)$; inflow $u_{\mathrm{in}}=4y(1-y)$; $p=0$ at the
outlet and no-slip $\mathbf{u}=0$ on the walls; $Re=100$.

\section{Wave Equation (hyperbolic)}
\paragraph{Wave1d-C.} $u_{tt}-4u_{xx}=0$, $(x,t)\in[0,1]\times[0,1]$,
$u(x,0)=\sin(\pi x)+\frac12\sin(4\pi x)$, $u_t(x,0)=0$, $u(0,t)=u(1,t)=0$; exact solution
\[
  u(x,t)=\sin(\pi x)\cos(2\pi t)+\tfrac12\sin(4\pi x)\cos(8\pi t).
\]

\paragraph{Wave2d-CG (heterogeneous medium).} $\nabla^2u-\frac{1}{c(\mathbf{x})}\frac{\partial^2u}{\partial t^2}=0$,
$(x,y)\in[-1,1]^2$, $t\in[0,5]$, $u(\mathbf{x},0)=\exp(-\norm{\mathbf{x}-\bm\mu}^2/(2\sigma^2))$, $u_t(\mathbf{x},0)=0$, Neumann
condition $\partial u/\partial\mathbf{n}=0$; $\bm\mu=(-0.5,0)$, $\sigma=0.3$; $c(\mathbf{x})$ a Gaussian random field.

\paragraph{Wave2d-MS (multi-scale).} Stated equation $u_{tt}-(u_{xx}+a^2u_{yy})=0$,
$(x,y,t)\in[0,1]^2\times[0,100]$, with reference solution
\[
  u=c_1\sin(m_1\pi x)\sinh(n_1\pi y)\cos(p_1\pi t)+c_2\sinh(m_2\pi x)\sin(n_2\pi y)\cos(p_2\pi t),
\]
$a=\sqrt2$, $m_1=1$, $m_2=3$, $n_1=1$, $n_2=2$, $p_1=p_2=1$, $c_1=c_2=1$.

\paragraph{The reference solution of Wave2d-MS is a Klein--Gordon solution.}
For the first term, $u_1=\sin(\pi x)\sinh(\pi y)\cos(\pi t)$, we have $\partial_{tt}u_1=-\pi^2u_1$,
$\partial_{xx}u_1=-\pi^2u_1$ and $\partial_{yy}u_1=\pi^2u_1$, so
\[
  \partial_{tt}u_1-(\partial_{xx}u_1+2\partial_{yy}u_1)=(-\pi^2+\pi^2-2\pi^2)u_1=-2\pi^2u_1.
\]
For the second term, $u_2=\sinh(3\pi x)\sin(2\pi y)\cos(\pi t)$, we have $\partial_{tt}u_2=-\pi^2u_2$,
$\partial_{xx}u_2=9\pi^2u_2$ and $\partial_{yy}u_2=-4\pi^2u_2$, so
\[
  \partial_{tt}u_2-(\partial_{xx}u_2+2\partial_{yy}u_2)=(-\pi^2-9\pi^2+8\pi^2)u_2=-2\pi^2u_2.
\]
Therefore, by linearity, $u=u_1+u_2$ satisfies
\begin{equation}
  u_{tt}-(u_{xx}+2u_{yy})+2\pi^2u=0,
  \label{eq:wave2dms-kg}
\end{equation}
which is a Klein--Gordon equation with mass $m^2=2\pi^2$, and not the stated wave equation. (The
computation was also checked symbolically.) This would explain why all the methods of
Table~\ref{tab:benchmark} report L2RE $\approx1$ in this task, and it is one more motivation for studying the
Klein--Gordon equation (Section~\ref{sec:klein-gordon}).

\end{hyphenrules}
